% TOP_PROBLEM: {"problemId": "problem.kannan-lovasz-simonovits-conjecture-on-isoperimetry-of-log-concave-measures", "canonicalTitle": "Kannan-Lovasz-Simonovits conjecture on isoperimetry of log-concave measures", "releaseRank": 98, "primaryDomain": "applied_computational_mathematics", "primaryDomainLabel": "Applied & computational mathematics", "displayStatus": "open", "releaseStatus": "open"}
% !TeX program = lualatex
% ATTEMPT_STATUS: unresolved
\documentclass[11pt]{article}
\usepackage[margin=25mm]{geometry}
\usepackage{fontspec}
\setmainfont{Latin Modern Roman}
\usepackage{amsmath,amssymb,mathtools}
\usepackage{unicode-math}
\setmathfont{Latin Modern Math}
\usepackage{xurl,hyperref}
\hypersetup{colorlinks=true,urlcolor=blue}
\setcounter{secnumdepth}{0}
\setlength{\parindent}{0pt}
\setlength{\parskip}{5pt}
\setlength{\emergencystretch}{3em}
\begin{document}
\section*{\texorpdfstring{Kannan-Lovasz-Simonovits conjecture on isoperimetry of log-concave measures}{Kannan-Lovasz-Simonovits conjecture on isoperimetry of log-concave measures}}\label{98}
\subsection{Definitions and mathematical statement}
For a log-concave probability density $e^{-V}$ on its affine support, $V$ is convex. Define the outer boundary measure by $\mu^+(B)=\liminf_{\varepsilon\downarrow0}(\mu(B_\varepsilon)-\mu(B))/\varepsilon$, and
\[
 h(\mu)=\inf_{0<\mu(B)<1}\frac{\mu^+(B)}{\min(\mu(B),1-\mu(B))}.
\]
The halfspace form asks for $h(\mu)\ge c\,h_{\rm half}(\mu)$, where the latter infimum is restricted to halfspaces and $c>0$ is universal. The dimension-free covariance formulation is $h(\mu)\ge c/\sqrt{\|\operatorname{Cov}(\mu)\|_{\rm op}}$. Work on the affine support for degenerate measures; point masses require the corresponding trivial convention.
\par\smallskip\textit{Formulation and concept references:} \href{https://arxiv.org/pdf/2603.29571}{[S1]}.

\subsection{Short English statement}
For log-concave distributions, are halfspaces always within a universal factor of the worst possible isoperimetric bottleneck?
\subsection{Research attempt}

\paragraph{Scope, process, and checked context.}
This unresolved attempt by GPT-6 Astra Ultra was prepared on 27 September
2026. The constant sought is independent of dimension and of the
log-concave probability measure. All distances and covariance operators
are Euclidean on the affine support; a point mass is separated as the
trivial zero-dimensional case. The attack first proves the
one-dimensional statement with explicit constants, then tries to pass
from one-dimensional control to arbitrary cuts.

The primary localization source is Kannan--Lovasz--Simonovits
\href{https://www.renyi.hu/~miki/KannanLovSimIso.pdf}{[E1]}; Bobkov
\href{https://www-users.cse.umn.edu/~bobkov/papers/2007_LNM_Isop.log-concave.pdf}{[E2]}
explains related isoperimetric reductions. The recent manuscript [S2]
states a logarithmic, dimension-dependent bound and a result for
quadratic forms. The thin-shell result in [S4] controls a particular
radial variance. Neither statement is substituted for the requested
uniform bound on all measurable cuts. The calculations below use no
claimed general solution of KLS.

\paragraph{A quantile calculation.}
Let $f$ be a nondegenerate log-concave density on an interval of
$\mathbb R$, with distribution function $F$ and median $m$.
Put $c=f(m)>0$ and, for $0<u<1$, define
\[
                         J(u)=f(F^{-1}(u)).
\]
The function $J$ is concave. Indeed, on the interior of the support
$\log f$ is concave and locally Lipschitz, and the chain rule gives
$J'(u)=(\log f)'(F^{-1}(u))$ almost everywhere. The right side is
nonincreasing; the locally absolutely continuous function $J$ is
therefore concave. This also covers nonsmooth log densities by
one-sided derivatives. Endpoint limits are nonnegative, and concavity
with $J(1/2)=c$ yields
\begin{equation}\label{r98:profile}
  J(u)\ge2c\min(u,1-u),\qquad
                           \sup_{0<u<1}J(u)\le2c.
\end{equation}
For the upper bound, when $u<1/2$ apply concavity between $u$ and a
point tending to $1$ to get $c\ge J(u)/(2(1-u))$; the other half is
symmetric. These are not assertions that the density itself is
symmetric or that its maximum occurs at the median.

\paragraph{Lemma 98.1: the exact one-dimensional Cheeger constant.}
For this measure,
\begin{equation}\label{r98:exact}
                              h(\mu)=2f(m).
\end{equation}
Here is a proof that includes arbitrary measurable cuts, not just
half-lines. For any locally Lipschitz $g$, integration along intervals
from $m$ and Tonelli's theorem give
\begin{align*}
 \int |g(x)-g(m)|f(x)\,dx
 &\le \int_{t<m}|g'(t)|F(t)\,dt
       +\int_{t>m}|g'(t)|(1-F(t))\,dt\\
 &\le \frac1{2c}\int |g'(t)|f(t)\,dt.
\end{align*}
The last inequality is exactly the lower bound in
\eqref{r98:profile}, on the appropriate side of the median.

For a measurable $B$, if its closure has strictly larger measure,
$\mu(B_\varepsilon)-\mu(B)$ is bounded below by a positive constant
and $\mu^+(B)=\infty$. Otherwise replace it by its closure without
changing its measure or its open enlargements, and use
\[
                 g_\varepsilon(x)=
                 \max\{0,1-\operatorname{dist}(x,B)/\varepsilon\}.
\]
Its derivative vanishes almost everywhere off
$B_\varepsilon\setminus B$, and has magnitude at most
$1/\varepsilon$ there. Moreover $g_\varepsilon$ tends to $1_B$ in
$L^1(\mu)$. For every $a\in[0,1]$,
$\int|1_B-a|\,d\mu\ge\min(\mu(B),1-\mu(B))$; applying this with
$a=g_\varepsilon(m)$ in the preceding functional inequality and
taking a liminf proves
\[
                  \mu^+(B)\ge2c\min(\mu(B),1-\mu(B)).
\]
The half-line $(-\infty,m]$ has mass $1/2$ and outer boundary measure
$f(m)=c$, so it attains the bound and proves \eqref{r98:exact}.
The reasoning takes place on the support and also covers compactly
supported densities.\hfill$\square$

\paragraph{Explicit comparison with variance.}
Write $\sigma^2=\operatorname{Var}(X)$. Chebyshev gives mass at least
$3/4$ to the interval of length $4\sigma$ centred at the mean.
Consequently $\|f\|_\infty\ge3/(16\sigma)$. By
\eqref{r98:profile}, $\|f\|_\infty\le2c$, giving
$h(\mu)\sigma\ge3/16$.

For the reverse bound, on $x\ge m$ the same profile estimate says
$f(x)\ge2c(1-F(x))$. Solving the resulting differential inequality
for the survival function, and doing the analogous calculation on
the left, yields
\[
 \mathbb P(X\ge m+t)\le\tfrac12e^{-2ct},\qquad
 \mathbb P(X\le m-t)\le\tfrac12e^{-2ct}\quad(t\ge0).
\]
Thus
\[
 \sigma^2\le\mathbb E(X-m)^2
   =\int_0^\infty2t\,\mathbb P(|X-m|>t)\,dt
   \le\frac1{2c^2}.
\]
Combining these inequalities gives the fully explicit, nonoptimal
two-sided comparison
\begin{equation}\label{r98:variance}
                  \frac3{16}\le h(\mu)\sigma\le\sqrt2.
\end{equation}
No dimension parameter has appeared because this calculation is
one-dimensional.

For the uniform measure on $[-a,a]$, the formula gives $h=1/a$ and
$\sigma=a/\sqrt3$. For the density
$f(x)=e^{-|x|/b}/(2b)$ it gives $h=1/b$ and
$\sigma=\sqrt2b$, attaining the upper constant above. For the
one-sided exponential $f(x)=e^{-x}$ on $[0,\infty)$, the median is
$\log2$ and $h=1$, even though the density maximum is at zero.
These examples check symmetry, support boundaries and normalization.

\paragraph{What this proves about halfspaces in every dimension.}
Let $\Sigma$ be positive definite on the affine support and put
$s=\sqrt{\|\Sigma\|_{\rm op}}$. The marginal theorem of Prekopa
\href{https://acta.bibl.u-szeged.hu/14411/1/math_034_335-343.pdf}{[E3]} for
log-concave densities makes the law of $\langle\theta,X\rangle$
log-concave for every unit vector $\theta$. Its variance is
$\sigma_\theta^2=\theta^{\mathsf T}\Sigma\theta$.
For a halfspace perpendicular to $\theta$, Euclidean enlargement
translates its threshold by exactly $\varepsilon$, so its boundary
ratio is the corresponding one-dimensional ratio. By
\eqref{r98:exact} its infimum over thresholds is attained at the
median. Therefore \eqref{r98:variance} gives
\begin{equation}\label{r98:half}
              \frac3{16s}\le h_{\rm half}(\mu)\le\frac{\sqrt2}{s}.
\end{equation}
The lower bound uses $\sigma_\theta\le s$ in every direction; the
upper bound chooses an eigenvector for the largest eigenvalue.

This verifies the constant-factor relationship between the two
catalogue formulations. If $h\ge c_0h_{\rm half}$, then
$h\ge3c_0/(16s)$. Conversely, $h\ge c_1/s$ and the upper bound in
\eqref{r98:half} give $h\ge(c_1/\sqrt2)h_{\rm half}$. What has not
been proved is either premise for arbitrary sets. The elementary
inequality from taking an infimum over more sets goes in the other
direction, $h\le h_{\rm half}$.

\paragraph{A precise obstruction to the naive localization step.}
It is tempting to apply the one-dimensional lemma to every needle
arising in localization and declare its variance bounded by isotropy
of the original measure. That declaration is false even for a cube.
Let $\mu$ be uniform on $[-\sqrt3,\sqrt3]^n$. Its covariance is $I$.
Along the diagonal unit vector $\theta=(1,\ldots,1)/\sqrt n$, the
central chord consists of $t\theta$ with
$|t|\le\sqrt{3n}$. The uniform probability measure on this chord
has
\[
          \operatorname{Var}(t)=n,\qquad
                             h=\frac1{\sqrt{3n}}.
\]
It is a legitimate one-dimensional log-concave measure on a segment
inside the same body. By contrast, the \emph{projection} of the
original cube measure onto $\theta$ has variance one. A restriction
to a chord and a marginal are different probability measures.

This does not say that a particular localization construction must
select that chord, or that the cube violates KLS. It disproves the
proposed justification that every possible needle inherits the
ambient covariance bound. A localization proof must preserve
appropriate moment constraints or otherwise control the bad
needles; the one-dimensional theorem alone supplies neither.

There is also a basic scaling check. Under $x\mapsto ax+b$ with
$a>0$, neighbourhood radii scale by $a$, so $h$ scales by $a^{-1}$
and $s$ by $a$. For a general invertible linear map $A$, if its
smallest and largest singular values are $s_{\min},s_{\max}$, the
inclusions between transformed neighbourhoods give
\[
 \frac{h(\mu)}{s_{\max}}
 \le h(A_*\mu)\le\frac{h(\mu)}{s_{\min}}.
\]
Thus whitening is useful but does not leave the Euclidean boundary
constant unchanged. All appearances of isotropy must account for
that metric change.

\paragraph{Verification and first gap.}
The exact one-dimensional formula was proved through a functional
inequality and distance cutoffs, including nonsmooth cuts. The
variance bounds use only concavity and elementary tail integration;
the examples and diagonal-chord calculation are exact. The marginal
step imports preservation of log-concavity under linear projection,
but it makes no claim about conditional needle variances. These are
classical ingredients and transparent deductions, with no novelty
claim or numerical evidence of a general solution.

The first missing inequality is a dimension-independent comparison
between arbitrary-cut boundary ratios and the halfspace ratios in
\eqref{r98:half}. Next targets are a localization procedure with
controlled needle variances, an all-functions functional inequality
strong enough to test distance cutoffs, and a search for isotropic
log-concave examples whose nonlinear cuts outperform halfspaces by
an unbounded factor. Quadratic or radial tests alone do not cover
those distance functions, and the diagonal chord is not a
counterexample to the original measure's inequality.

\paragraph{Final status.}
\textbf{UNRESOLVED.} The one-dimensional formula and variance bounds,
the halfspace/covariance comparison, and the failure of a naive
needle argument are verified. General isoperimetric control in
arbitrary dimension remains unproved.

\subsection{Sources}
\begin{itemize}
\item[S1]Afonso S. Bandeira, Daniil Dmitriev, Kevin Lucca, Petar Nizić-Nikolac, and Almut Rödder, "Randomstrasse101: Open Problems of 2025," arXiv:2603.29571, 2026.
\url{https://arxiv.org/pdf/2603.29571}.
\textit{Formulation source.}
\item[S2]The KLS constant is O(log\textasciicircum{}\{1/4\} n).
\url{https://arxiv.org/abs/2607.24164}.
\textit{Further reference; not a proof endorsement.}
\item[S3]Spectral and Isoperimetric Bounds on Flat Tori.
\url{https://arxiv.org/html/2608.13052v1}.
\textit{Further reference; not a proof endorsement.}
\item[S4]Digesting the proof of the sharp thin-shell inequality.
\url{https://www.weizmann.ac.il/math/klartag/sites/math.klartag/files/uploads/sharp_thin.pdf}.
\textit{Further reference; not a proof endorsement.}
\item[E1]R. Kannan, L. Lovasz and M. Simonovits, Isoperimetric Problems for Convex Bodies and a Localization Lemma, Discrete and Computational Geometry 13 (1995), 541--559.
\url{https://www.renyi.hu/~miki/KannanLovSimIso.pdf}.
\textit{Primary source for the localization framework; consulted 27 September 2026.}
\item[E2]S. G. Bobkov, On Isoperimetric Constants for Log-Concave Probability Distributions, Geometric Aspects of Functional Analysis, Lecture Notes in Mathematics 1910 (2007), 81--88.
\url{https://www-users.cse.umn.edu/~bobkov/papers/2007_LNM_Isop.log-concave.pdf}.
\textit{Primary context for one-dimensional bounds and their localization limits; consulted 27 September 2026.}
\item[E3]A. Prekopa, On logarithmic concave measures and functions, Acta Scientiarum Mathematicarum 34 (1973), 335--343.
\url{https://acta.bibl.u-szeged.hu/14411/1/math_034_335-343.pdf}.
\textit{Primary source for preservation of log-concavity under marginalization; consulted 27 September 2026.}
\end{itemize}
\end{document}
